\onehalfspacing The A-/Ā-distinction—a distinction between dependency types like raising-to-subject and wh-movement—has been a mainstay of modern syntactic theorizing since \textcite{chomsky1981LecturesGovernmentBinding, safir2019DistinctionEpiphenomenon}. Yet, despite widespread acceptance of the empirical correlates of this distinction, its theoretical underpinning remains a matter of debate. In this talk, I offer a novel perspective on understanding the A-/Ā-distinction by examining resumptive A-dependencies--- that is, A-dependencies terminating in a resumptive pronoun---in Iraqi Arabic. A-resumption challenges traditional positional approaches to the A-/Ā-distinction that distinguish just two types of positions \parencite[i.e. A vs. Ā; see e.g.][]{chomsky1981LecturesGovernmentBinding, mahajan1990AbarDistinctionMovementa} because A-resumption displays a mixture of A- and Ā-properties. Moreover, A-resumption challenges approaches which derive A- and Ā-properties from the mechanics of movement \parencite[see e.g.][]{vanurk2015UniformSyntaxPhrasal} because A-resumption does not involve movement. Instead, I develop a featural analysis of A-resumption and the A-/Ā-distinction according to which (i) most differences between A- and Ā-dependencies are linked to the category of the head of the dependency—D and Q, respectively—and (ii) differences between movement-derived and base-generated dependencies stem from a difference between the featural triggers for external and internal Merge \parencite[building on][]{hewettTypesResumptiveADependencies}. A major upshot of my investigation is the discovery that certain properties previously analyzed as characterizing A- and Ā-dependencies are better understood as characterizing A- and Ā-movement, thereby elucidating our understanding of dependency types cross-linguistically.